\documentclass[10pt,a4paper]{amsart}
\usepackage{amsmath,amssymb,amsthm,graphicx}
\usepackage{verbatim}
\newenvironment{reference}{\comment}{\endcomment}
%\newenvironment{reference}{}{}
\newenvironment{slogan}{\comment}{\endcomment}
\newenvironment{history}{\comment}{\endcomment}

% For commutative diagrams we use Xy-pic
\usepackage[all]{xy}

% We use 2cell for 2-commutative diagrams.
\xyoption{2cell}
% Theorem environments.
%
\theoremstyle{plain}
\newtheorem{theorem}[subsection]{Theorem}
\newtheorem{proposition}[subsection]{Proposition}
\newtheorem{lemma}[subsection]{Lemma}

\theoremstyle{definition}
\newtheorem{definition}[subsection]{Definition}
\newtheorem{example}[subsection]{Example}
\newtheorem{exercise}[subsection]{Exercise}
\newtheorem{situation}[subsection]{Situation}

\theoremstyle{remark}
\newtheorem{remark}[subsection]{Remark}
\newtheorem{remarks}[subsection]{Remarks}

\numberwithin{equation}{subsection}

% Macros
%
\def\lim{\mathop{\mathrm{lim}}\nolimits}
\def\colim{\mathop{\mathrm{colim}}\nolimits}
\def\Spec{\mathop{\mathrm{Spec}}}
\def\Hom{\mathop{\mathrm{Hom}}\nolimits}
\def\Ext{\mathop{\mathrm{Ext}}\nolimits}
\def\SheafHom{\mathop{\mathcal{H}\!\mathit{om}}\nolimits}
\def\SheafExt{\mathop{\mathcal{E}\!\mathit{xt}}\nolimits}
\def\Sch{\mathit{Sch}}
\def\Mor{\mathop{\mathrm{Mor}}\nolimits}
\def\Ob{\mathop{\mathrm{Ob}}\nolimits}
\def\Sh{\mathop{\mathit{Sh}}\nolimits}
\def\NL{\mathop{N\!L}\nolimits}
\def\CH{\mathop{\mathrm{CH}}\nolimits}
\def\proetale{{pro\text{-}\acute{e}tale}}
\def\etale{{\acute{e}tale}}
\def\QCoh{\mathit{QCoh}}
\def\Ker{\mathop{\mathrm{Ker}}}
\def\Im{\mathop{\mathrm{Im}}}
\def\Coker{\mathop{\mathrm{Coker}}}
\def\Coim{\mathop{\mathrm{Coim}}}

% Boxtimes
%
\DeclareMathSymbol{\boxtimes}{\mathbin}{AMSa}{"02}

%
% Macros for moduli stacks/spaces
%
\def\QCohstack{\mathcal{QC}\!\mathit{oh}}
\def\Cohstack{\mathcal{C}\!\mathit{oh}}
\def\Spacesstack{\mathcal{S}\!\mathit{paces}}
\def\Quotfunctor{\mathrm{Quot}}
\def\Hilbfunctor{\mathrm{Hilb}}
\def\Curvesstack{\mathcal{C}\!\mathit{urves}}
\def\Polarizedstack{\mathcal{P}\!\mathit{olarized}}
\def\Complexesstack{\mathcal{C}\!\mathit{omplexes}}
% \Pic is the operator that assigns to X its picard group, usage \Pic(X)
% \Picardstack_{X/B} denotes the Picard stack of X over B
% \Picardfunctor_{X/B} denotes the Picard functor of X over B
\def\Pic{\mathop{\mathrm{Pic}}\nolimits}
\def\Picardstack{\mathcal{P}\!\mathit{ic}}
\def\Picardfunctor{\mathrm{Pic}}
\def\Deformationcategory{\mathcal{D}\!\mathit{ef}}
\usepackage{cleveref}
\crefname{originallemma}{original lemma}{original lemmas}
\Crefname{originallemma}{Original lemma}{Original lemmas}
\makeatletter
\newlabel{lemma-domination-universally-injective}{{0AUM}{source}{}{sourcepointer-1}{}}
\newlabel{lemma-domination-universally-injective@cref}{{[originallemma][0][]0AUM}{[1][0][]source}{}{}{}}
\makeatother
\crefname{originallemma}{original lemma}{original lemmas}
\Crefname{originallemma}{Original lemma}{Original lemmas}
\makeatletter
\newlabel{proposition-ML-characterization}{{059E}{source}{}{sourcepointer-2}{}}
\newlabel{proposition-ML-characterization@cref}{{[originallemma][0][]059E}{[1][0][]source}{}{}{}}
\makeatother
\title{The Mittag-Leffler property descends along a faithfully flat ring map}
\author{The Stacks Project Authors}
\date{stacks-project:algebra.tex:local-snapshot-20261012; GNU FDL 1.2 or later}
\begin{document}
\maketitle
\noindent\textit{Editorial note.} All delivered mathematical text is the author's own, copied from the pinned Stacks Project chapter file algebra.tex; no statement or proof step was rewritten, reordered or removed.\par\medskip
\noindent\textit{Editorial note.} Preamble provenance, as required for this source: the frozen snapshot stages the chapter's own preamble file beside it (preamble.tex, 217 lines, 8116 bytes, the same shared Stacks-wide preamble the published chow and etale-cohomology cards use); the first line of the pinned chapter file is \textbackslash{}input\{preamble\}, so no preamble text lives inside the chapter itself. The card takes that file's lines 8-18 (the verbatim package and the three \textbackslash{}comment-based environments reference/slogan/history at lines 9-12, section 13.2, together with the Xy-pic setup) and its lines 154 to the runtime end of the file (the native theorem declarations and mathematical macros). Not one character was changed, no package was added, no macro was defined, and the excerpt serves only as preamble, never as a span coordinate.\par\medskip
\noindent\textit{Editorial note.} The same-chapter pointers are located at their \textbackslash{}label lines inside the pinned chapter file; the cross-chapter pointers are located in the project's own tag list staged in the snapshot (sources/stacks-chow/tags.txt), at the verbatim lines carrying their canonical tags.\par\medskip
\noindent\textit{Editorial note.} The counted statement contains a reference environment recording the origin of the statement. That environment is defined by the chapter's own preamble as a comment environment (lines 9-12, inside this card's preamble spans), so it is not typeset; the text is kept verbatim inside the span.\par\medskip
\noindent\textit{Editorial note.} The counted proof refers to three same-chapter statements by label and contains two Xy-pic pushout squares; the statements are recorded as navigation stubs at their \textbackslash{}label lines and the Xy-pic package is provided by the preamble spans.\par\medskip
\noindent\textit{Allowed standard prerequisites (editorial).}\begin{itemize}
\item The characterisation of Mittag-Leffler modules by domination. Invoked in the counted proof at lines 22570-22572. Reduces the claim to finding an index over which a comparison map dominates another.
\item Domination tested by universal injectivity of a pushout. Invoked in the counted proof at lines 22582-22584. Turns the domination requirement into universal injectivity of a map out of the pushout.
\end{itemize}\medskip
\ifdefined\hypertarget\hypertarget{sourcepointer-1}{}\fi
\noindent\textit{Source reference.} Original lemma 0AUM: sources/stacks-algebra/pinned-original-source.tex, source lines 21153--21153. Navigation pointer for \textbackslash{}ref\{lemma-domination-universally-injective\} used inside the counted span (line 22584); the argument that consumes it is retained in the span.\par\medskip
\ifdefined\hypertarget\hypertarget{sourcepointer-2}{}\fi
\noindent\textit{Source reference.} Original lemma 059E: sources/stacks-algebra/pinned-original-source.tex, source lines 21217--21217. Navigation pointer for \textbackslash{}ref\{proposition-ML-characterization\} used inside the counted span (line 22570); the argument that consumes it is retained in the span.\par\medskip
\begin{lemma}
\label{lemma-ffdescent-ML}
\begin{reference}
Email from Juan Pablo Acosta Lopez dated 12/20/14.
\end{reference}
Let $R \to S$ be a faithfully flat ring map. Let $M$ be an $R$-module. If the
$S$-module $M \otimes_R S$ is Mittag-Leffler, then $M$ is Mittag-Leffler.
\end{lemma}

\begin{proof}
Write $M = \colim_{i\in I} M_i$ as a directed colimit
of finitely presented $R$-modules $M_i$.
Using Proposition \ref{proposition-ML-characterization}, we see that we
have to prove that for each $i \in I$ there exists $i \leq j$, $j\in I$
such that $M_i\rightarrow M_j$ dominates $M_i\rightarrow M$.

\medskip\noindent
Take $N$ the pushout
$$
\xymatrix{
M_i  \ar[r] \ar[d] & M_j \ar[d] \\
M \ar[r] & N
}
$$
Then the lemma is equivalent to the existence of $j$ such that
$M_j\rightarrow N$ is universally injective, see
Lemma \ref{lemma-domination-universally-injective}.
Observe that the tensorization by $S$
$$
\xymatrix{
M_i\otimes_R S  \ar[r] \ar[d] & M_j\otimes_R S \ar[d] \\
M\otimes_R S \ar[r] & N\otimes_R S
}
$$
Is a pushout diagram.  So because
$M \otimes_R S = \colim_{i\in I} M_i \otimes_R S$
expresses $M\otimes_R S$ as a colimit of $S$-modules of
finite presentation, and $M\otimes_R S$ is Mittag-Leffler,
there exists $j \geq i$ such that $M_j\otimes_R S\rightarrow N\otimes_R S$
is universally injective. So using that $R\rightarrow S$ is faithfully flat
we conclude that $M_j\rightarrow N$ is universally injective too.
\end{proof}

\end{document}
